% Capítulo gerado a partir de ponto_flutuante.tex (arquivos usados neste capítulo)
\def\arqBase{MudancaDeBase.ggb}
\def\arqSisF{SistemaPontoFlutuante_reta.ggb}
\def\arqMoito{Maquina_PontoFlutuante_8bits.ggb}
\def\arqMdezesseis{Maquina_PontoFlutuante_16bits.ggb}
\def\arqMtrintadois{Maquina_PontoFlutuante_32bits.ggb}
\def\arqMsessentaquatro{Maquina_PontoFlutuante_64bits.ggb}
\providecommand{\bits}[1]{\texttt{#1}}
\chapter{Mudança de base e aritmética de ponto flutuante}\label{cap:pontoflutuante}

\begin{tcolorbox}[colback=gray!4,colframe=gray!50,title={Como usar os aplicativos e os textos deste capítulo},fonttitle=\bfseries]
Os quadros laranja \emph{Atividade no GeoGebra} trazem o link \textbf{[abrir o aplicativo]}.
São seis aplicativos. Os quatro primeiros substituem as planilhas de Excel das
\emph{máquinas} de 8, 16, 32 e 64 bits:
\begin{itemize}[nosep]
  \item máquinas de ponto flutuante: \ggbabrir{\arqMoito}{8 bits}, \ggbabrir{\arqMdezesseis}{16 bits},
        \ggbabrir{\arqMtrintadois}{32 bits (\emph{float})} e \ggbabrir{\arqMsessentaquatro}{64 bits (\emph{double})}.
        Cada bit é uma caixa de seleção; o aplicativo decodifica a palavra, codifica um número
        digitado (por arredondamento ou truncamento) e calcula os erros;
  \item \ggbabrir{\arqBase}{mudança de base}, com as divisões e as multiplicações sucessivas passo a passo;
  \item \ggbabrir{\arqSisF}{sistema $F(\beta,t,m,M)$ na reta}, que mostra todos os números de
        uma máquina pequena e o valor $fl(x)$ de qualquer $x$.
\end{itemize}

\end{tcolorbox}

% ==================================================================
\section{Por que estudar a representação dos números}

A resolução de um problema por métodos numéricos passa por várias fases: o problema real,
o levantamento de dados, a construção do modelo matemático, a escolha do método, a
implementação computacional e a análise dos resultados. Não é raro que o resultado final
fique distante do esperado, mesmo com todas as fases feitas corretamente. Os resultados
dependem também
\begin{itemize}[nosep]
  \item da precisão dos dados de entrada;
  \item da forma como esses dados são representados no computador;
  \item das operações numéricas efetuadas.
\end{itemize}

\begin{exemplo}[{: a área de um círculo}]
A área de um círculo de raio $100$\,m, calculada com $\pi\approx3{,}14$, $\pi\approx3{,}1416$ e
$\pi\approx3{,}141592654$, dá
\[
A=31\,400\ \text{m}^2,\qquad A=31\,416\ \text{m}^2,\qquad A=31\,415{,}9265\ \text{m}^2 .
\]
As diferenças vêm da aproximação usada para $\pi$, que é irracional. Nenhum computador obtém
a área ``exatamente''.
\end{exemplo}

\begin{exemplo}[{: um somatório}]
Calcule $S=\sum_{i=1}^{30\,000}x_i$ com $x_i=0{,}5$ e com $x_i=0{,}11$. As respostas exatas são
$15\,000$ e $3\,300$. Somando num programa:

\begin{center}
\begin{tabular}{lcc}
\toprule
 & $x_i=0{,}5$ & $x_i=0{,}11$\\
\midrule
precisão simples (32 bits) & $15\,000$ & $3\,300{,}985$\\
precisão dupla (64 bits)   & $15\,000$ & $3\,300{,}000\,000\,000\,628\,5$\\
\bottomrule
\end{tabular}
\end{center}

Com $0{,}5$ o resultado é exato; com $0{,}11$, não. A explicação está neste capítulo:
$0{,}5=(0{,}1)_2$ tem representação binária finita, mas $0{,}11$ não tem. Cada parcela já entra
com um pequeno erro, e cada soma acrescenta outro.
\end{exemplo}



Há três fontes de erro num cálculo numérico:
\begin{enumerate}[nosep]
  \item \textbf{incerteza dos dados}: os dados vêm de medidas com instrumentos de precisão finita;
  \item \textbf{erros de arredondamento}: vêm da forma finita de representar números na máquina;
  \item \textbf{erros de truncamento}: aparecem quando um processo infinito (uma série, um limite,
        uma integral) é substituído por um processo finito.
\end{enumerate}
Este capítulo trata do segundo tipo.

% ==================================================================
\section{Sistemas de numeração posicionais}

A \emph{base} de um sistema de numeração é a quantidade de símbolos usados para escrever
os números.

\begin{center}
\begin{tabular}{lll}
\toprule
Base & Nome & Símbolos\\
\midrule
2  & binário      & 0, 1\\
8  & octal        & 0, 1, 2, 3, 4, 5, 6, 7\\
10 & decimal      & 0, 1, \dots, 9\\
16 & hexadecimal  & 0, 1, \dots, 9, A, B, C, D, E, F\\
\bottomrule
\end{tabular}
\end{center}

A base 10 é a do nosso dia a dia. A base 2 é natural nos circuitos lógicos: 0 e 1 correspondem a
dois níveis de tensão (\emph{low} e \emph{high}), a falso e verdadeiro, a desligado e ligado.
As bases 8 e 16 servem para escrever números binários de forma compacta. Quando a base não é
10, ela vem indicada no índice: $(1010001011)_2$, $(453234)_8$, $(23A0D)_{16}$.

\begin{definicao}[notação posicional]
Num sistema posicional de base $\beta$, o número
$(d_{p-1}d_{p-2}\dots d_1d_0\,,\,d_{-1}\dots d_{-n})_\beta$, com $0\le d_i\le\beta-1$, vale
\[
D=\sum_{i=-n}^{p-1}d_i\,\beta^{\,i}
 =d_{p-1}\beta^{p-1}+\dots+d_1\beta+d_0+d_{-1}\beta^{-1}+\dots+d_{-n}\beta^{-n},
\]
em que $p$ é o número de dígitos à esquerda da vírgula e $n$ o número de dígitos à direita.
O valor de cada símbolo depende da sua posição.
\end{definicao}

Em base 2, cada dígito é um \emph{bit}. O bit mais à esquerda é o MSB
(\emph{most significant bit}) e o mais à direita é o LSB (\emph{least significant bit}).

\begin{center}
\small
\begin{tabular}{cccc|cccc}
\toprule
hex & dec & bin & & hex & dec & bin\\
\midrule
0 & 0 & 0000 && 8 & 8  & 1000\\
1 & 1 & 0001 && 9 & 9  & 1001\\
2 & 2 & 0010 && A & 10 & 1010\\
3 & 3 & 0011 && B & 11 & 1011\\
4 & 4 & 0100 && C & 12 & 1100\\
5 & 5 & 0101 && D & 13 & 1101\\
6 & 6 & 0110 && E & 14 & 1110\\
7 & 7 & 0111 && F & 15 & 1111\\
\bottomrule
\end{tabular}
\end{center}

% ==================================================================
\section{Da base \texorpdfstring{$\beta$}{beta} para a base 10}

Basta somar os dígitos multiplicados pelos pesos $\beta^{\,i}$.

\begin{exemplo}
\begin{align*}
(11011)_2  &=1\cdot2^4+1\cdot2^3+0\cdot2^2+1\cdot2^1+1\cdot2^0=16+8+2+1=27,\\
(372)_8    &=3\cdot8^2+7\cdot8+2=192+56+2=250,\\
(24{,}6)_8 &=2\cdot8+4+6\cdot8^{-1}=16+4+0{,}75=20{,}75,\\
(356)_{16} &=3\cdot16^2+5\cdot16+6=768+80+6=854,\\
(2AF)_{16} &=2\cdot16^2+10\cdot16+15=512+160+15=687,\\
(1000{,}011)_2&=8+0\cdot2^{-1}+1\cdot2^{-2}+1\cdot2^{-3}=8+0{,}25+0{,}125=8{,}375.
\end{align*}
\end{exemplo}

% ==================================================================
\section{Da base 10 para a base \texorpdfstring{$\beta$}{beta}}

O número é separado em parte inteira e parte fracionária, que são convertidas por processos
diferentes.

\subsection{Parte inteira: divisões sucessivas}

Seja $N$ inteiro com $(N)_{10}=(c_k\dots c_2c_1c_0)_\beta$. Então
\[
N=c_k\beta^k+\dots+c_2\beta^2+c_1\beta+c_0=\beta\,\underbrace{(c_k\beta^{k-1}+\dots+c_2\beta+c_1)}_{q_1}+c_0 .
\]
Logo $c_0$ é o \emph{resto} da divisão de $N$ por $\beta$ e $q_1$ é o quociente. Repetindo
o raciocínio com $q_1$, obtemos $c_1$, e assim por diante, até o quociente ser zero.

\begin{tcolorbox}[colback=azul!3,colframe=azul!60,title={Algoritmo: inteiro decimal $\to$ base $\beta$},fonttitle=\bfseries]
\begin{description}[nosep,font=\normalfont\bfseries]
  \item[Passo 0:] $k=0$, $N_0=N$.
  \item[Passo 1:] obtenha $q_k$ e $r_k$ tais que $N_k=\beta\,q_k+r_k$, $0\le r_k<\beta$; faça $c_k=r_k$.
  \item[Passo 2:] se $q_k=0$, pare. Caso contrário, faça $N_{k+1}=q_k$, $k=k+1$ e volte ao passo 1.
\end{description}
Os restos, lidos \textbf{de baixo para cima} (do último para o primeiro), formam o número na base $\beta$.
\end{tcolorbox}

\begin{exemplo}
\[
\begin{array}{rcl}
25&=&2\cdot12+\mathbf{1}\\
12&=&2\cdot6+\mathbf{0}\\
6&=&2\cdot3+\mathbf{0}\\
3&=&2\cdot1+\mathbf{1}\\
1&=&2\cdot0+\mathbf{1}
\end{array}
\ \Rightarrow\ 25=(11001)_2,
\qquad
\begin{array}{rcl}
2834&=&8\cdot354+\mathbf{2}\\
354&=&8\cdot44+\mathbf{2}\\
44&=&8\cdot5+\mathbf{4}\\
5&=&8\cdot0+\mathbf{5}
\end{array}
\ \Rightarrow\ 2834=(5422)_8 .
\]
Da mesma forma, $37=(100101)_2$ e $2834=16\cdot177+2$, $177=16\cdot11+1$, $11=16\cdot0+11$,
logo $2834=(B12)_{16}$.
\end{exemplo}

\subsection{Parte fracionária: multiplicações sucessivas}

Seja $0<r<1$ com $(r)_{10}=(0{,}d_1d_2d_3\dots)_\beta$. Multiplicando por $\beta$:
\[
\beta\,r=d_1+\underbrace{(0{,}d_2d_3\dots)_\beta}_{\text{parte fracionária}} .
\]
A parte inteira de $\beta r$ é $d_1$. Repetindo com a parte fracionária, obtemos $d_2$, $d_3$,\dots

\begin{tcolorbox}[colback=azul!3,colframe=azul!60,title={Algoritmo: fração decimal $\to$ base $\beta$},fonttitle=\bfseries]
\begin{description}[nosep,font=\normalfont\bfseries]
  \item[Passo 0:] $r_1=r$, $k=1$.
  \item[Passo 1:] calcule $\beta r_k$; faça $d_k=\lfloor\beta r_k\rfloor$ (parte inteira).
  \item[Passo 2:] faça $r_{k+1}=\beta r_k-d_k$. Se $r_{k+1}=0$, pare.
  \item[Passo 3:] faça $k=k+1$ e volte ao passo 1.
\end{description}
As partes inteiras, lidas \textbf{de cima para baixo}, formam $0{,}d_1d_2d_3\dots$ Se o processo não
terminar, paramos depois do número de dígitos desejado.
\end{tcolorbox}

\begin{exemplo}[{: $8{,}375$ e $4{,}8$ em binário}]
Parte inteira: $8=(1000)_2$ e $4=(100)_2$. Partes fracionárias:
\[
\begin{array}{rcll}
0{,}375\times2&=&\mathbf{0}{,}75\\
0{,}75\times2&=&\mathbf{1}{,}5\\
0{,}5\times2&=&\mathbf{1}{,}0&\text{(fim)}
\end{array}
\qquad\qquad
\begin{array}{rcll}
0{,}8\times2&=&\mathbf{1}{,}6\\
0{,}6\times2&=&\mathbf{1}{,}2\\
0{,}2\times2&=&\mathbf{0}{,}4\\
0{,}4\times2&=&\mathbf{0}{,}8&\text{(volta ao início)}
\end{array}
\]
Logo $8{,}375=(1000{,}011)_2$, uma expansão finita, e
$4{,}8=(100{,}\overline{1100})_2=(100{,}110011001100\dots)_2$, uma expansão infinita e periódica.
\end{exemplo}

\begin{exemplo}[{: $3{,}38$ e $57{,}3$ (Aula 01)}]
Para $0{,}38$, as multiplicações dão $0{,}76$; $1{,}52$; $1{,}04$; $0{,}08$; $0{,}16$; $0{,}32$; $0{,}64$; $1{,}28$;\dots,
e então $0{,}38\approx(0{,}01100001)_2=0{,}37890625$ com 8 dígitos. Portanto
$3{,}38\approx(11{,}01100001)_2$. Quanto mais dígitos, melhor a aproximação.

Para $57{,}3$: $57=(111001)_2$ e $0{,}3=(0{,}0\overline{1001})_2$, logo
$57{,}3=(111001{,}0100110011001\dots)_2$.
\end{exemplo}

\begin{observacao}
Uma fração $p/q$ irredutível tem expansão finita na base $\beta$ se, e somente se, todo fator primo
de $q$ divide $\beta$. Em base 2, só têm expansão finita os racionais com denominador potência de 2:
$1/8=(0{,}001)_2$, $1/16=(0{,}0001)_2$, $5{,}625=(101{,}101)_2$. Já $1/3$, $1/5$ e $1/10$ são periódicos:
\[
\tfrac13=(0{,}\overline{01})_2,\qquad \tfrac15=(0{,}\overline{0011})_2,\qquad
\tfrac1{10}=(0{,}0\overline{0011})_2 .
\]
Um número com expansão decimal finita, como $0{,}1$, pode ter expansão binária infinita.
\end{observacao}

\subsection{Entre as bases 2, 8 e 16}

Como $8=2^3$ e $16=2^4$, basta agrupar os bits de 3 em 3 (octal) ou de 4 em 4 (hexadecimal) a partir da
vírgula:
\[
(101101{,}101)_2=(\underbrace{101}_{5}\,\underbrace{101}_{5}{,}\underbrace{101}_{5})_2=(55{,}5)_8,
\qquad
(0010\,1101{,}1010)_2=(2D{,}A)_{16} .
\]
É por isso que as palavras de máquina são escritas em hexadecimal.

\begin{atividade}{mudança de base passo a passo}{\arqBase}
Digite um número decimal (até 9 casas decimais), escolha a base $\beta$ (de 2 a 16) e o número de
dígitos fracionários. O aplicativo mostra:
\begin{itemize}[nosep]
  \item à esquerda, as divisões sucessivas da parte inteira, com os restos em vermelho;
  \item no centro, as multiplicações sucessivas da parte fracionária;
  \item à direita, o resultado, a verificação pela notação posicional, o erro cometido ao parar e a
        conferência com o comando \texttt{ToBase} do GeoGebra;
  \item no quadro de baixo, a conversão de um inteiro de uma base qualquer para decimal
        (\texttt{FromBase}).
\end{itemize}
\textbf{Tarefas.} (a) Reproduza os exemplos $25$, $37$, $2834$ (bases 8 e 16), $8{,}375$, $4{,}8$,
$3{,}38$ e $57{,}3$.
(b) Converta $0{,}1$ para a base 2 com 4, 8, 16 e 24 dígitos e anote o erro em cada caso.
Quantos dígitos binários são necessários para que o erro fique abaixo de $10^{-6}$?
(c) Ache três números com expansão finita em base 2 e infinita em base 3.
\end{atividade}

% ==================================================================
\section{Representação em ponto flutuante}

Um número real $x\neq0$ pode ser escrito em \emph{ponto flutuante normalizado} na base $\beta$:
\[
x=\pm\,(0{,}d_1d_2\dots d_t)_\beta\times\beta^{\,e},\qquad d_1\neq0,\qquad -m\le e\le M,
\]
em que $0{,}d_1\dots d_t$ é a \emph{mantissa}, $t$ é o número de dígitos da mantissa e $e$ é o
\emph{expoente}, limitado por $-m$ e $M$. Esse conjunto de números é indicado por $F(\beta,t,m,M)$.
O zero é representado à parte, com mantissa nula.

Como a mantissa e o expoente têm tamanho finito:
\begin{itemize}[nosep]
  \item $F$ é um conjunto \emph{finito}, com $2(\beta-1)\beta^{t-1}(M+m+1)+1$ elementos;
  \item o menor positivo é $(0{,}10\dots0)_\beta\times\beta^{-m}=\beta^{-m-1}$, e o maior é
        $(0{,}(\beta{-}1)\dots(\beta{-}1))_\beta\times\beta^{M}=(1-\beta^{-t})\beta^{M}$;
  \item se $0<|x|<\beta^{-m-1}$, ocorre \textbf{underflow}; se $|x|>(1-\beta^{-t})\beta^{M}$,
        ocorre \textbf{overflow};
  \item os números de $F$ \textbf{não são igualmente espaçados}: entre $\beta^{e-1}$ e $\beta^{e}$ o
        espaçamento é $\beta^{e-t}$, que cresce com $e$.
\end{itemize}

\begin{exemplo}[{: $F(10,3,5,5)$}]
Com $t=3$ dígitos decimais e $-5\le e\le5$, o menor positivo é $0{,}100\times10^{-5}=10^{-6}$ e o maior é
$0{,}999\times10^{5}=99\,900$. Então
$x=235{,}89=0{,}23589\times10^{3}$ fica $0{,}236\times10^3$ (arredondado) ou $0{,}235\times10^3$
(truncado); $x=0{,}345\times10^{-7}$ dá underflow; $x=0{,}875\times10^{9}$ dá overflow.
\end{exemplo}

\begin{atividade}{o sistema $F(\beta,t,m,M)$ na reta}{\arqSisF}
Os controles deslizantes escolhem $\beta$, $t$, $m$ e $M$; os pontos azuis do eixo são \emph{todos}
os números do sistema. Arraste o ponto $P$: o ponto vermelho é $fl(x)$, por arredondamento ou,
desmarcando a caixa, por truncamento. O quadro mostra a quantidade de números, o maior, o menor
positivo, os erros absoluto e relativo e os avisos de underflow e overflow.

\textbf{Tarefas.} (a) Em $F(2,3,1,2)$, confira que há 33 números, que o maior é $3{,}5$ e o menor
positivo é $0{,}25$. (b) Amplie a região perto do zero (roda do mouse) e observe o ``buraco''
entre $0$ e $0{,}25$. (c) Compare o espaçamento entre $0{,}5$ e $1$ com o espaçamento entre $2$ e $3{,}5$.
(d) Mova $P$ por vários pontos e verifique que $|ER|$ nunca passa de $\tfrac12\beta^{1-t}$ com
arredondamento, enquanto $|EA|$ cresce com $|x|$.
\end{atividade}

% ==================================================================
\section{O padrão IEEE 754}

As mais diversas representações em ponto flutuante já foram usadas. O padrão IEEE~754, criado em 1985
e revisto em 2008 e 2019, é hoje implementado na grande maioria dos processadores. Em base 2, a
mantissa normalizada começa sempre com o bit 1, que por isso não precisa ser guardado (\emph{bit
escondido}). O padrão usa a forma $1{,}f$ e guarda três campos:
\[
\underbrace{\ s\ }_{\text{sinal}}\ \Big|\ \underbrace{\ \ E\ \ }_{\text{expoente com bias}}\ \Big|\ \underbrace{\qquad f\qquad}_{\text{fração}}
\qquad\qquad
x=(-1)^s\times(1{,}f)_2\times2^{E-\text{bias}} .
\]

\begin{itemize}[nosep]
  \item $s=0$ para positivo e $s=1$ para negativo;
  \item $E$ é um inteiro sem sinal; o expoente verdadeiro é $E-\text{bias}$, com
        $\text{bias}=2^{k-1}-1$ para um campo de $k$ bits (assim não é preciso guardar expoentes negativos);
  \item $f=(0{,}b_1b_2\dots b_p)_2$ é a fração; a mantissa é $1{,}f$, com $p+1$ bits de precisão.
\end{itemize}

\begin{center}
\begin{tabular}{lccccc}
\toprule
formato & total & expoente & fração & bias & faixa aproximada\\
\midrule
meia (\emph{binary16})            & 16  & 5  & 10  & 15    & $6{,}1\times10^{-5}$ a $65\,504$\\
simples (\texttt{float})          & 32  & 8  & 23  & 127   & $1{,}2\times10^{-38}$ a $3{,}4\times10^{38}$\\
dupla (\texttt{double})           & 64  & 11 & 52  & 1023  & $2{,}2\times10^{-308}$ a $1{,}8\times10^{308}$\\
estendida (Intel)                 & 80  & 15 & 63  & 16383 & $3{,}4\times10^{-4932}$ a $1{,}2\times10^{4932}$\\
quádrupla (\emph{binary128})      & 128 & 15 & 112 & 16383 & idem\\
\bottomrule
\end{tabular}
\end{center}

\subsection{Valores especiais e números subnormais}

Os valores extremos do campo $E$ são reservados:

\begin{center}
\begin{tabular}{lll}
\toprule
campo $E$ & fração $f$ & significado\\
\midrule
$00\dots0$ & $f=0$     & $\pm0$\\
$00\dots0$ & $f\neq0$  & subnormal: $x=(-1)^s\times(0{,}f)_2\times2^{\,1-\text{bias}}$\\
$1\le E\le 2^k-2$ & qualquer & normalizado: $x=(-1)^s\times(1{,}f)_2\times2^{E-\text{bias}}$\\
$11\dots1$ & $f=0$     & $\pm\infty$ (overflow, divisão por zero)\\
$11\dots1$ & $f\neq0$  & NaN, \emph{not a number} (por exemplo, $0/0$ ou $\sqrt{-1}$)\\
\bottomrule
\end{tabular}
\end{center}

Os \textbf{subnormais} (ou desnormalizados) preenchem o intervalo entre $0$ e o menor normalizado
$2^{1-\text{bias}}$ com números igualmente espaçados. O underflow fica ``gradual'', em vez de saltar
direto para zero.

\begin{exemplo}[{: $15\,213$ em precisão simples}]
$15\,213=(11101101101101)_2=(1{,}1101101101101)_2\times2^{13}$. Então
\begin{itemize}[nosep]
  \item $s=0$;
  \item $E=13+127=140=(10001100)_2$;
  \item $f=\bits{11011011011010000000000}$ (os 13 bits depois da vírgula, completados com zeros até 23).
\end{itemize}
A palavra é \bits{0\,10001100\,11011011011010000000000}, ou $466DB400_{16}$.
\end{exemplo}

\begin{exemplo}[{: $178{,}125$ em precisão dupla}]
$178{,}125=128+32+16+2+0{,}125=(10110010{,}001)_2=(1{,}0110010001)_2\times2^{7}$. Então
$s=0$, $E=7+1023=1030=(10000000110)_2$ e $f=\bits{0110010001}0\dots0$ (52 bits).
A palavra é $4066440000000000_{16}$.
\end{exemplo}

\begin{exemplo}[{: $0{,}2$ em todos os formatos (as planilhas das máquinas)}]
$0{,}2=(0{,}0\overline{0011})_2=(1{,}\overline{1001})_2\times2^{-3}$. O expoente é $-3$ e a fração
$1001\,1001\,1001\dots$ precisa ser cortada. Arredondando para o mais próximo:

\begin{center}
\small
\begin{tabular}{lccc}
\toprule
formato & palavra (hex) & valor guardado & erro relativo\\
\midrule
8 bits (bias 3)   & $0D$          & $0{,}203125$ & $-1{,}6\times10^{-2}$\\
16 bits           & $3266$        & $0{,}199951171875$ & $\phantom{-}2{,}4\times10^{-4}$\\
32 bits           & $3E4CCCCD$    & $0{,}20000000298023224$ & $-1{,}5\times10^{-8}$\\
64 bits           & $3FC999999999999A$ & $0{,}20000000000000001110$ & $-5{,}6\times10^{-17}$\\
\bottomrule
\end{tabular}
\end{center}

Na máquina de 8 bits, o expoente $-3$ é menor que o mínimo $1-3=-2$, e $0{,}2$ cai na faixa
subnormal: $0{,}2\approx(0{,}1101)_2\times2^{-2}=13/64$. Nas planilhas de Excel o campo $E=0$ era tratado
como normalizado. No padrão IEEE (e nos aplicativos) ele é reservado para o zero e os subnormais.
\end{exemplo}



% ==================================================================
\section{As máquinas de 8, 16, 32 e 64 bits}

As planilhas usadas em aula montavam um número bit a bit e calculavam
$(-1)^s(1+\text{fração})\,2^{E-\text{bias}}$. Os aplicativos a seguir fazem o mesmo e mais:
\begin{itemize}[nosep]
  \item cada bit é uma caixa de seleção (marcada $=1$), agrupada em sinal, expoente e fração, com o
        peso de cada bit indicado acima da caixa;
  \item a palavra aparece em binário e em hexadecimal;
  \item a \emph{decodificação} mostra $s$, $E$, $E-\text{bias}$, a fração, o tipo do número (zero,
        subnormal, normalizado, infinito ou NaN), o valor com 15 algarismos significativos e o
        espaçamento local (ulp, \emph{unit in the last place});
  \item na \emph{codificação}, você digita o número exato e escolhe \textbf{arredondamento} (para o mais
        próximo, empate para o par, como no padrão) ou \textbf{truncamento} (corte, em direção ao zero).
        O aplicativo marca os bits e calcula os erros absoluto e relativo;
  \item o quadro \emph{Faixa} mostra o maior número, o menor normalizado e o menor subnormal.
\end{itemize}

\subsection{A máquina de 8 bits}

O formato didático de 8 bits de W.~T.~Verts tem 1 bit de sinal, 3 de expoente
($\text{bias}=3$) e 4 de fração. Ele segue todas as regras do IEEE~754 e é pequeno o bastante para
listar os seus 256 padrões:
\begin{itemize}[nosep]
  \item $E=0$: zero e subnormais $(0{,}f)_2\times2^{-2}$, de $1/64=0{,}015625$ a $15/64$, com passo $1/64$;
  \item $1\le E\le6$: normalizados $(1{,}f)_2\times2^{E-3}$, de $0{,}25$ a $(1{,}1111)_2\times2^3=15{,}5$;
  \item $E=7$: $\pm\infty$ e NaN.
\end{itemize}
São $2\times(16+6\times16)=224$ padrões finitos (com $+0$ e $-0$). Entre $1$ e $2$ o passo é $1/16$;
entre $8$ e $15{,}5$ é $1/2$.

\begin{center}
\small
\begin{tabular}{cccl}
\toprule
$E$ & expoente & intervalo & passo\\
\midrule
0 (sub.) & $-2$ & $[0,\ 0{,}234375]$ & $1/64$\\
1 & $-2$ & $[0{,}25,\ 0{,}484375]$ & $1/64$\\
2 & $-1$ & $[0{,}5,\ 0{,}96875]$ & $1/32$\\
3 & $\phantom{-}0$ & $[1,\ 1{,}9375]$ & $1/16$\\
4 & $\phantom{-}1$ & $[2,\ 3{,}875]$ & $1/8$\\
5 & $\phantom{-}2$ & $[4,\ 7{,}75]$ & $1/4$\\
6 & $\phantom{-}3$ & $[8,\ 15{,}5]$ & $1/2$\\
\bottomrule
\end{tabular}
\end{center}

\begin{atividade}{máquina de 8 bits}{\arqMoito}
O aplicativo abre com $0{,}2$ arredondado: \bits{0 000 1101}, um subnormal que vale $0{,}203125$.

\textbf{Tarefas.} (a) Clique no bit de expoente de peso $2^0$ e veja o número passar a
$(1{,}1101)_2\times2^{-2}=0{,}453125$. (b) Monte o maior número ($15{,}5$), o menor normalizado ($0{,}25$),
o menor subnormal ($1/64$), $+\infty$ e um NaN. (c) Digite $x=2{,}7$ e codifique por arredondamento e por
truncamento: os resultados são $2{,}75$ e $2{,}625$. (d) Codifique $15{,}75$ por arredondamento (overflow,
$+\infty$) e por truncamento (fica o maior finito, $15{,}5$). (e) Codifique $0{,}007$, que é menor que
metade do menor subnormal: o resultado é zero (underflow).
\end{atividade}



\subsection{As máquinas de 16, 32 e 64 bits}

\begin{atividade}{máquina de 16 bits (meia precisão)}{\arqMdezesseis}
Formato $1+5+10$, $\text{bias}=15$. O aplicativo abre com $0{,}2$ arredondado
($3266_{16}$, $0{,}199951171875$).

\textbf{Tarefas.} (a) Monte $1$, $2$, $-1{,}5$ e o maior número finito ($65\,504$). (b) Codifique
$65\,519$ e $65\,520$ por arredondamento: o primeiro vira $65\,504$ e o segundo, $+\infty$.
(c) Codifique $2049$: o resultado é $2048$, porque com 11 bits de mantissa os inteiros acima de
$2^{11}$ não são todos representáveis. (d) Compare o erro relativo de $0{,}2$ com $\varepsilon_{\text{máq}}=2^{-10}$.
\end{atividade}

\begin{atividade}{máquina de 32 bits (precisão simples)}{\arqMtrintadois}
Formato $1+8+23$, $\text{bias}=127$, o \texttt{float} das linguagens de programação.

\textbf{Tarefas.} (a) Codifique $15\,213$ e confira a palavra $466DB400_{16}$. (b) Codifique $0{,}1$ e
$0{,}2$ e anote os erros relativos; compare com $2^{-24}\approx5{,}96\times10^{-8}$.
(c) Codifique $16\,777\,217=2^{24}+1$ e explique o resultado. (d) Codifique $10^{-40}$: o número é
subnormal e o erro relativo cresce, porque sobram menos bits significativos.
\end{atividade}

\begin{atividade}{máquina de 64 bits (precisão dupla)}{\arqMsessentaquatro}
Formato $1+11+52$, $\text{bias}=1023$, o \texttt{double}. Os 52 bits da fração estão em duas linhas
(use as barras de rolagem se a tela for pequena). Em hexadecimal, a palavra tem 16 dígitos e a fração 13.

\textbf{Tarefas.} (a) Codifique $178{,}125$ e confira $4066440000000000_{16}$. (b) Codifique $0{,}1$:
a fração é $999999999999A_{16}$; o último dígito ``A'' (em vez de 9) vem do arredondamento para cima.
(c) Compare a palavra de $0{,}1$ com a de $0{,}2$: só muda o expoente. Por quê?
(d) O GeoGebra calcula em precisão dupla, então o número digitado já entra como \texttt{double}: aqui o
erro aparece como $0$. Use o ulp e $\varepsilon_{\text{máq}}=2^{-52}$ para estimar o erro, ou faça a mesma conta na
máquina de 32 bits.
\end{atividade}

% ==================================================================
\section{Arredondamento e truncamento}

Considere um sistema que opera em aritmética de ponto flutuante de $t$ dígitos na base 10, e escreva
\[
x=f_x\times10^{e}+g_x\times10^{e-t},\qquad 0{,}1\le f_x<1,\quad 0\le g_x<1 .
\]
Por exemplo, com $t=4$ e $x=234{,}57$: $x=0{,}2345\times10^{3}+0{,}7\times10^{-1}$, logo $f_x=0{,}2345$ e $g_x=0{,}7$.
A parcela $g_x\times10^{e-t}$ não cabe na mantissa. Há dois critérios:

\begin{description}[nosep]
  \item[Truncamento:] $g_x$ é desprezado e $\bar x=f_x\times10^{e}$. Então
  \[
  |EA_x|=|g_x|\times10^{e-t}<10^{e-t},\qquad
  |ER_x|=\frac{|g_x|\times10^{e-t}}{|f_x|\times10^{e}}<\frac{10^{e-t}}{0{,}1\times10^{e}}=10^{-t+1} .
  \]
  \item[Arredondamento:] se $|g_x|<\tfrac12$, despreza-se $g_x$; caso contrário, soma-se 1 ao último dígito de
  $f_x$. Então
  \[
  |EA_x|\le\tfrac12\times10^{e-t},\qquad |ER_x|<\tfrac12\times10^{-t+1} .
  \]
\end{description}

O arredondamento comete erros menores, mas exige mais operações; por isso algumas máquinas antigas
truncavam. O IEEE~754 adota como padrão o \textbf{arredondamento para o mais próximo, com empate para o
par} (o último bit fica 0). Ele também oferece arredondamento em direção a $0$ (truncamento), a $+\infty$ e a
$-\infty$, úteis na análise por intervalos.

Na base 2 com $p$ bits de fração, os mesmos argumentos dão o \textbf{épsilon da máquina}
$\varepsilon_{\text{máq}}=2^{-p}$, a distância de $1$ ao número seguinte da máquina, e
\[
\frac{|x-fl(x)|}{|x|}\le\tfrac12\varepsilon_{\text{máq}}=2^{-(p+1)}\ \ \text{(arredondamento)},
\qquad
\frac{|x-fl(x)|}{|x|}<\varepsilon_{\text{máq}}\ \ \text{(truncamento)},
\]
para $x$ na faixa dos normalizados. Em precisão dupla, $\varepsilon_{\text{máq}}=2^{-52}\approx2{,}2\times10^{-16}$,
ou seja, cerca de 16 algarismos decimais significativos. Em precisão simples,
$\varepsilon_{\text{máq}}=2^{-23}\approx1{,}2\times10^{-7}$, cerca de 7 algarismos.

% ==================================================================
\section{Erros absoluto e relativo}

\begin{definicao}
Seja $\bar x$ uma aproximação de $x$. O \emph{erro absoluto} é $EA_x=x-\bar x$ e o \emph{erro relativo} é
\[
ER_x=\frac{EA_x}{\bar x}=\frac{x-\bar x}{\bar x}
\]
(na prática, divide-se pelo valor aproximado, que é o conhecido; quando $x$ é conhecido, como nos aplicativos,
pode-se dividir por $x$).
\end{definicao}

Em geral só $\bar x$ é conhecido, e o que se obtém é um limitante do erro. Se
$\pi\in(3{,}14;\,3{,}15)$ e tomamos $\bar\pi$ nesse intervalo, então $|EA_\pi|<0{,}01$.

\begin{exemplo}[{: o erro relativo compara melhor}]
Seja $\bar x=2112{,}9$ com $|EA_x|<0{,}1$, isto é, $x\in(2112{,}8;\,2113)$, e $\bar y=5{,}3$ com
$|EA_y|<0{,}1$, isto é, $y\in(5{,}2;\,5{,}4)$. Os erros absolutos têm o mesmo limitante, mas
\[
|ER_x|<\frac{0{,}1}{2112{,}9}\approx4{,}7\times10^{-5},\qquad |ER_y|<\frac{0{,}1}{5{,}3}\approx0{,}02 .
\]
O número $x$ está representado com precisão muito maior que $y$.
\end{exemplo}

% ==================================================================
\section{Operações em ponto flutuante e propagação de erros}

Mesmo quando as parcelas são exatas, o resultado de uma operação precisa ser normalizado e arredondado (ou
truncado) para $t$ dígitos. O resultado aproximado $\overline{OP}$ é o que fica na memória.

\begin{exemplo}[{: adição e multiplicação com $t=4$}]
Sejam $x=0{,}937\times10^{4}$ e $y=0{,}1272\times10^{2}$.

\textbf{Adição.} É preciso alinhar as vírgulas: a mantissa do número de menor expoente é deslocada para a
direita tantas casas quanto a diferença dos expoentes:
\[
x+y=(0{,}937+0{,}001272)\times10^{4}=0{,}938272\times10^{4} .
\]
Com $t=4$: $\overline{x+y}=0{,}9383\times10^{4}$ (arredondamento) ou $0{,}9382\times10^{4}$ (truncamento).

\textbf{Multiplicação.}
$xy=(0{,}937\times0{,}1272)\times10^{6}=0{,}1191864\times10^{6}$, logo
$\overline{xy}=0{,}1192\times10^{6}$ (arredondamento) ou $0{,}1191\times10^{6}$ (truncamento).
\end{exemplo}

Supondo as parcelas representadas exatamente, o erro relativo de uma operação é
$|ER_{OP}|<10^{-t+1}$ no truncamento e $|ER_{OP}|<\tfrac12\times10^{-t+1}$ no arredondamento.

\subsection{Fórmulas de propagação}

Sejam $x=\bar x+EA_x$ e $y=\bar y+EA_y$. Desprezando os produtos de erros, que são muito pequenos:

\begin{center}
\renewcommand{\arraystretch}{2}
\begin{tabular}{lll}
\toprule
operação & erro absoluto & erro relativo\\
\midrule
$x+y$ & $EA_{x+y}=EA_x+EA_y$ & $ER_{x+y}=ER_x\dfrac{\bar x}{\bar x+\bar y}+ER_y\dfrac{\bar y}{\bar x+\bar y}$\\
$x-y$ & $EA_{x-y}=EA_x-EA_y$ & $ER_{x-y}=ER_x\dfrac{\bar x}{\bar x-\bar y}-ER_y\dfrac{\bar y}{\bar x-\bar y}$\\
$xy$  & $EA_{xy}\approx\bar x\,EA_y+\bar y\,EA_x$ & $ER_{xy}\approx ER_x+ER_y$\\
$x/y$ & $EA_{x/y}\approx\dfrac{\bar y\,EA_x-\bar x\,EA_y}{\bar y^{\,2}}$ & $ER_{x/y}\approx ER_x-ER_y$\\
\bottomrule
\end{tabular}
\end{center}

Na divisão, usa-se $\dfrac1{1+EA_y/\bar y}=1-\dfrac{EA_y}{\bar y}+\Big(\dfrac{EA_y}{\bar y}\Big)^2-\cdots$ e
desprezam-se as potências maiores que 1. A cada operação soma-se ainda o erro de arredondamento do próprio
resultado, limitado por $\tfrac12\varepsilon_{\text{máq}}$.

\subsection{Cancelamento}

Na subtração, os fatores $\bar x/(\bar x-\bar y)$ e $\bar y/(\bar x-\bar y)$ ficam enormes quando
$\bar x\approx\bar y$. Os algarismos iniciais iguais se cancelam e sobram os finais, justamente os
contaminados por erro. É o \textbf{cancelamento}.

\begin{exemplo}[{: reescrever para evitar o cancelamento}]
Com $t=6$ dígitos decimais, calcule $\sqrt{x+1}-\sqrt{x}$ para $x=10^{8}$:
$\sqrt{100\,000\,001}=10\,000{,}000\,05\approx0{,}100000\times10^{5}$ e $\sqrt{10^8}=0{,}100000\times10^{5}$, cuja
diferença é $0$. Racionalizando,
\[
\sqrt{x+1}-\sqrt{x}=\frac{1}{\sqrt{x+1}+\sqrt{x}}\approx\frac1{0{,}200000\times10^{5}}=0{,}500000\times10^{-4},
\]
que é o valor correto ($4{,}9999999875\times10^{-5}$) com 6 algarismos.
\end{exemplo}

\begin{exemplo}[{: derivada numérica de $\sin x$ em $x=1$}]
A derivada é aproximada por $\dfrac{\sin(1+h)-\sin 1}{h}$; o valor exato é $\cos1\approx0{,}5403023$. Em precisão dupla:

\begin{center}
\small
\begin{tabular}{ccc}
\toprule
$h$ & aproximação & erro\\
\midrule
$10^{-1}$  & $0{,}4973637525$ & $4{,}3\times10^{-2}$\\
$10^{-4}$  & $0{,}5402602314$ & $4{,}2\times10^{-5}$\\
$10^{-8}$  & $0{,}5403023029$ & $3{,}0\times10^{-9}$\\
$10^{-12}$ & $0{,}5403455461$ & $4{,}3\times10^{-5}$\\
$10^{-15}$ & $0{,}5551115123$ & $1{,}5\times10^{-2}$\\
$10^{-16}$ & $0$ & $5{,}4\times10^{-1}$\\
\bottomrule
\end{tabular}
\end{center}

Até $h\approx10^{-8}\approx\sqrt{\varepsilon_{\text{máq}}}$, o erro cai junto com $h$ (erro de truncamento da
fórmula). Depois, $\sin(1+h)$ e $\sin1$ ficam tão próximos que o cancelamento domina. Com
$h=10^{-16}$, $1+h$ é arredondado para $1$ e a diferença é zero. A diferença central
$\frac{\sin(1+h)-\sin(1-h)}{2h}$ atinge erro de $10^{-11}$ com $h\approx10^{-5}$.
\end{exemplo}

\begin{observacao}
Outros cuidados que vêm da aritmética finita:
\begin{itemize}[nosep]
  \item a adição não é associativa: $(x+y)+z$ pode ser diferente de $x+(y+z)$;
  \item somar um número muito pequeno a um muito grande pode não mudar nada: em precisão dupla,
        $1+2^{-53}=1$;
  \item não teste igualdade de números reais (\texttt{if x == 0.3}); use uma tolerância,
        $|x-0{,}3|<\varepsilon$. Em Python, \texttt{0.1 + 0.2} mostra \texttt{0.30000000000000004}.
\end{itemize}
\end{observacao}



% ==================================================================
\section{Exercícios}

\begin{enumerate}[label=\arabic*.]
  \item Converta para a base 10: (a) $(101101)_2$; (b) $(0{,}1011)_2$; (c) $(157)_8$; (d) $(3A{,}8)_{16}$.
  \item Converta $45{,}625$ para as bases 2, 8 e 16.
  \item Mostre que $0{,}1=(0{,}0\overline{0011})_2$. Quantos bits seriam necessários para guardar
        $0{,}1$ exatamente?
  \item Represente $13{,}25$ em precisão simples (32 bits) e dê a palavra em hexadecimal.
  \item Decodifique a palavra de 32 bits \bits{0 10000001 01000000000000000000000}.
  \item Na máquina de 8 bits ($1+3+4$, bias 3), decodifique:
        (a) \bits{0 101 0110}; (b) \bits{1 001 1000}; (c) \bits{0 000 0011}; (d) \bits{0 111 0000}.
  \item Na máquina de 8 bits, codifique $x=2{,}7$ por arredondamento e por truncamento e calcule os erros
        absoluto e relativo.
  \item Em $F(2,3,1,2)$, quantos números existem? Qual o maior e qual o menor positivo?
  \item Com $t=4$ e arredondamento, calcule $(x+y)-x$ para $x=0{,}3491\times10^{4}$ e $y=0{,}2345\times10^{0}$.
        Compare com $y$.
  \item Explique por que, em precisão dupla, $1+2^{-53}=1$, mas $1+2^{-52}\neq1$.
  \item Calcule $\sqrt{x+1}-\sqrt{x}$ para $x=10^8$ com $t=6$ dígitos, diretamente e após racionalizar.
\end{enumerate}

\subsection*{Respostas}
\begin{enumerate}[label=\arabic*.,nosep]
  \item (a) $45$; (b) $0{,}6875$; (c) $111$; (d) $58{,}5$.
  \item $(101101{,}101)_2=(55{,}5)_8=(2D{,}A)_{16}$.
  \item As multiplicações dão a sequência $0{,}2\to0{,}4\to0{,}8\to0{,}6\to0{,}2$, que se repete. Nenhuma
        quantidade finita de bits basta.
  \item $13{,}25=(1101{,}01)_2=(1{,}10101)_2\times2^3$: $s=0$, $E=130=(10000010)_2$,
        $f=\bits{10101}0\dots0$; palavra $41540000_{16}$.
  \item $(1{,}01)_2\times2^{129-127}=1{,}25\times4=5$.
  \item (a) $(1{,}0110)_2\times2^{2}=5{,}5$; (b) $-(1{,}1)_2\times2^{-2}=-0{,}375$; (c) subnormal
        $3/64=0{,}046875$; (d) $+\infty$.
  \item $2{,}7=(1{,}0101\,1001\dots)_2\times2^1$. Arredondamento: \bits{0 100 0110}$=2{,}75$, $EA=-0{,}05$,
        $ER\approx-1{,}85\times10^{-2}$. Truncamento: \bits{0 100 0101}$=2{,}625$, $EA=0{,}075$, $ER\approx2{,}78\times10^{-2}$.
  \item $2\cdot1\cdot4\cdot4+1=33$ números; maior $(0{,}111)_2\times2^2=3{,}5$; menor positivo $(0{,}100)_2\times2^{-1}=0{,}25$.
  \item $x+y=3491{,}2345\to0{,}3491\times10^{4}$, logo $(x+y)-x=0$, embora $y\neq0$: o $y$ foi ``absorvido''.
  \item O número seguinte a $1$ em precisão dupla é $1+2^{-52}$. O valor $1+2^{-53}$ está exatamente no meio e,
        pelo empate para o par, é arredondado para $1$.
  \item Diretamente: $0$. Racionalizando: $0{,}500000\times10^{-4}$ (exato: $4{,}9999999875\times10^{-5}$).
\end{enumerate}

\subsection*{Referências}
\begin{itemize}[nosep]
  \item M.~A.~G.~Ruggiero e V.~L.~R.~Lopes, \emph{Cálculo Numérico: aspectos teóricos e computacionais}, 2.ª ed., Makron Books.
  \item M.~L.~Overton, \emph{Numerical Computing with IEEE Floating Point Arithmetic}, SIAM, 2001.
  \item IEEE Std 754-2019, \emph{IEEE Standard for Floating-Point Arithmetic}.
\end{itemize}
